%! Function Model Selection and Construction — Unit 1, Lesson 1.7 — Guided Notes
\documentclass[10pt]{article}
\usepackage{precalculus-article}
\usepackage{precalculus-boxes}
\newcommand{\TallMath}[1]{$\displaystyle #1\rule[-1.4em]{0pt}{3.2em}$}

\begin{document}

\pageheader{Unit 1, Lesson 1.7}{Guided Notes}

\vspace{0.12in}

\begin{objectivebox}
By the end of this lesson, I will be able to\ldots
\begin{enumerate}[itemsep=2pt]
  \item Decide which quantity in a situation is the \blank{2.2cm} --- the
        model's input --- and which quantities become \blank{2.2cm}.
  \item Choose a model type from a table, by taking \blank{1.8cm} and
        \blank{2cm} differences of \blank{2.4cm} spaced inputs.
  \item Choose a model type from the context alone, by asking how the
        \blank{1.8cm} moves as the \blank{2cm} grows.
  \item Build a model by writing the \blank{2.2cm}, solving it for the other
        quantities, and \blank{2.4cm} into the one I want.
  \item Use the model: evaluate it, find an average rate of change with
        \blank{1.6cm}, and state the \blank{2.2cm} the context allows.
\end{enumerate}
\end{objectivebox}

\vspace{0.1in}

\boxguard
\begin{vocabbox}
Fill in each term as we define it together.
\par\vspace{2pt}
\noindent\textbf{\textcolor{plum}{Function model:}}\\[1pt]
\writeline\\[3pt]
\noindent\textbf{\textcolor{plum}{Constraint:}}\\[1pt]
\writeline\\[3pt]
\noindent\textbf{\textcolor{plum}{First differences:}}\\[1pt]
\writeline\\[3pt]
\noindent\textbf{\textcolor{plum}{Second differences:}}\\[1pt]
\writeline\\[3pt]
\noindent\textbf{\textcolor{plum}{Domain restriction:}}\\[1pt]
\writeline\\[3pt]
\noindent\textbf{\textcolor{plum}{Model assumption:}}\\[1pt]
\writeline
\end{vocabbox}

\vspace{0.1in}

\boxguard[24]
\begin{hookbox}
The garden club has exactly \textbf{80 feet of fencing} for a rectangular pen.
One whole side runs along the barn wall, so that side needs no fence at all.

\smallskip
Two members walked the field and measured the \textbf{same five trial pens}.
Each one wrote down a different thing.

\smallskip
\renewcommand{\arraystretch}{1.3}
\begin{center}
\begin{tabular}{|c|c|c|c|c|c|}
\hline
\textbf{Ana:} width $w$ (ft) & 5 & 10 & 15 & 20 & 25 \\
\hline
area $A$ (ft$^2$) & 350 & 600 & 750 & 800 & 750 \\
\hline
\end{tabular}

\smallskip
\begin{tabular}{|c|c|c|c|c|c|}
\hline
\textbf{Ben:} length $\ell$ (ft) & 30 & 40 & 50 & 60 & 70 \\
\hline
area $A$ (ft$^2$) & 750 & 800 & 750 & 600 & 350 \\
\hline
\end{tabular}
\end{center}

\smallskip
Same five pens. Same five areas. Two tables that look nothing alike. Which
member recorded the pen correctly?
\writelines{2}
\end{hookbox}

\vspace{0.1in}

\boxguard[30]
\begin{notesbox}{1. Pick the Question}
Every lesson so far handed you a function. Today you are handed a
\textbf{situation} --- and a situation is not a function yet.

\smallskip
\begin{center}
\begin{tikzpicture}[x=1cm, y=1cm]
  \fill[lilac] (0,0) rectangle (5,1.8);
  \draw[very thick, slate] (-0.5,1.8) -- (5.5,1.8);
  \node[slate, above, font=\small] at (2.5,1.85) {barn wall --- no fencing needed};
  \draw[very thick, plum] (0,1.8) -- (0,0) -- (5,0) -- (5,1.8);
  \node[plum, left,  font=\small] at (0,0.9) {$w$};
  \node[plum, right, font=\small] at (5,0.9) {$w$};
  \node[plum, below, font=\small] at (2.5,0) {$\ell$};
  \node[slate, font=\small\itshape] at (2.5,0.9) {area $A$};
\end{tikzpicture}
\end{center}

\smallskip
There are four numbers in this pen, and they are not all the same kind of
thing:

\smallskip
\renewcommand{\arraystretch}{1.35}
\begin{center}
\begin{tabularx}{0.94\linewidth}{@{} >{\bfseries}p{0.30\linewidth} X @{}}
\hline
80 feet of fencing & fixed --- it never changes, no matter which pen \\
\hline
the width $w$      & changes from pen to pen \\
\hline
the length $\ell$  & changes from pen to pen \\
\hline
the area $A$       & changes from pen to pen \\
\hline
\end{tabularx}
\end{center}

\medskip
A function takes \textbf{one} question and hands back \textbf{one} answer. So
before any algebra, you make two decisions.

\smallskip
\begin{center}
\begin{tikzpicture}[x=1cm, y=1cm]
  \draw[thick, plum, fill=lilac, rounded corners=2pt]
    (-3.6,-0.55) rectangle (-1.4,0.55);
  \node[font=\small] at (-2.5,0) {one quantity};
  \draw[thick, plum, fill=goldbg, rounded corners=2pt]
    (1.4,-0.55) rectangle (3.6,0.55);
  \node[font=\small] at (2.5,0) {one quantity};
  \draw[->, very thick, plum] (-1.3,0) -- (1.3,0);
  \node[above, font=\small] at (0,0.08) {the model};
  \node[below, font=\small\itshape, text=slate] at (-2.5,-0.62) {the question};
  \node[below, font=\small\itshape, text=slate] at (2.5,-0.62) {the answer};
\end{tikzpicture}
\end{center}

\medskip
Today the \textbf{answer} is settled --- the club cares about area. The
\textbf{question} is a choice, and it is yours to make.

\smallskip
\begin{center}
\begin{tcolorbox}[colback=lilac, colframe=plum, boxrule=0.8pt, arc=2mm,
    left=3mm, right=3mm, top=2mm, bottom=2mm, width=0.94\linewidth]
\textbf{A model starts with a decision.} Naming the question is the first move
in every modeling problem, and everything after it follows from that one
choice: which differences you may take, what the constraint has to do, which
inputs are allowed, and what the units say.
\end{tcolorbox}
\end{center}

\medskip
Back to the two notebooks.

\smallskip
\begin{itemize}[itemsep=4pt]
  \item Ana's question was the \blank{2.4cm}. \quad Ben's question was the
        \blank{2.4cm}.
  \item Both were writing about the \blank{2cm} five pens, so neither one is
        \blank{2cm}.
\end{itemize}

\medskip
And the 80? It is neither a question nor an answer. It is the
\textbf{constraint} --- the one sentence that stays true no matter which
question you pick. Write it now:

\begin{center}
\blank{1cm} $w \; + \;$ \blank{1cm} $= 80$
\end{center}
\end{notesbox}

\vspace{0.1in}

\boxguard[30]
\begin{notesbox}{2. When the Table Already Picked the Question}
A table has made the choice for you. The \textbf{top row is the question} and
the bottom row is the answer. All that is left is to ask how the answer moves.

\medskip
\textbf{Table A.} Take the first differences.

\smallskip
\renewcommand{\arraystretch}{1.4}
\begin{center}
\begin{tabular}{|c|c|c|c|c|c|}
\hline
$x$ & 0 & 1 & 2 & 3 & 4 \\
\hline
$y$ & 5 & 8 & 11 & 14 & 17 \\
\hline
1st difference & --- & \blank{1cm} & \blank{1cm} & \blank{1cm} & \blank{1cm} \\
\hline
\end{tabular}
\end{center}

\smallskip
The first differences are constant, so Table A is \blank{2.4cm}.

\medskip
\textbf{The requirement, and the reason for it.} A difference compares two
answers across \textit{one step of the question}. That comparison is only fair
when every step is the \blank{2.4cm} size --- so check the question row before
you subtract anything.

\smallskip
The warm-up's second table is the trap. Its differences were 4 and 8, which are
not constant --- so is $g$ not linear?

\smallskip
\renewcommand{\arraystretch}{1.4}
\begin{center}
\begin{tabular}{|c|c|c|c|}
\hline
$x$ & 1 & 2 & 4 \\
\hline
$g(x)$ & 3 & 7 & 15 \\
\hline
change in the question & --- & \blank{1cm} & \blank{1cm} \\
\hline
change in the answer & --- & 4 & 8 \\
\hline
answer change \textbf{per unit} & --- & \blank{1.1cm} & \blank{1.1cm} \\
\hline
\end{tabular}
\end{center}

\smallskip
Per unit of question the rate is the same, so $g$ \textit{is} \blank{2.2cm}
after all. The raw differences lied because the questions were not
\blank{3cm} spaced.

\medskip
\textbf{Ana's table} from the hook. Now the differences are fair --- her widths
step by 5 every time.

\smallskip
\renewcommand{\arraystretch}{1.4}
\begin{center}
\begin{tabular}{|c|c|c|c|c|c|}
\hline
$w$ (ft) & 5 & 10 & 15 & 20 & 25 \\
\hline
$A$ (ft$^2$) & 350 & 600 & 750 & 800 & 750 \\
\hline
1st difference & --- & \blank{1.1cm} & \blank{1.1cm} & \blank{1.1cm} & \blank{1.1cm} \\
\hline
2nd difference & --- & --- & \blank{1.1cm} & \blank{1.1cm} & \blank{1.1cm} \\
\hline
\end{tabular}
\end{center}

\smallskip
The first differences are \textbf{not} constant --- so take differences again.
The second differences are all \blank{1.6cm}, so Ana's data is
\blank{2.6cm}.

\medskip
\textbf{The decision tree.}
\begin{itemize}[itemsep=4pt]
  \item 1st differences constant $\Rightarrow$ \blank{2.4cm} \quad (degree \blank{0.8cm})
  \item 2nd differences constant $\Rightarrow$ \blank{2.4cm} \quad (degree \blank{0.8cm})
  \item 3rd differences constant $\Rightarrow$ \blank{2.4cm} \quad (degree \blank{0.8cm})
  \item $n$th differences constant $\Rightarrow$ a degree-\blank{0.8cm} polynomial
\end{itemize}

\medskip
\textbf{What it means.} Constant \textit{first} differences say the
\blank{2.4cm} of change repeats. Constant \textit{second} differences say the
\blank{2.4cm} of change is itself changing at a steady rate.

\medskip
\textbf{One warning about tables.} A degree-$n$ polynomial can be fit through
\blank{1.2cm} points with distinct inputs --- so any four points can be fit by
a cubic, and any five by a quartic. ``It fits the data'' is therefore not by
itself a reason to believe a model. The context gets a vote too, which is
Part 3.
\end{notesbox}

\vspace{0.1in}

\boxguard[28]
\begin{notesbox}{3. When Nobody Picked It: Reading the Context}
No table means you pick the question yourself --- and then the clue is in
\textbf{how the answer moves as the question grows}.

\smallskip
\renewcommand{\arraystretch}{1.5}
\begin{center}
\begin{tabularx}{0.94\linewidth}{@{} >{\raggedright\arraybackslash}p{0.52\linewidth} X @{}}
\hline
\textbf{What happens as the question grows} & \textbf{Model type} \\
\hline
The same amount is added each step & \blank{2.8cm} \\
\hline
\textbf{Two} lengths grow with it, and get multiplied & \blank{2.8cm} \\
\hline
\textbf{Three} lengths grow with it, and get multiplied & \blank{2.8cm} \\
\hline
The answer turns several times --- up, down, up & \blank{2.8cm} \\
\hline
The answer \textbf{shrinks}: the question is underneath & \blank{2.8cm} \\
\hline
The rule itself \textbf{changes at a threshold} & \blank{2.8cm} \\
\hline
\end{tabularx}
\end{center}

\medskip
\textbf{Careful --- the clue is not the noun.} ``Area'' does not mean quadratic
all by itself. What matters is how many lengths move when the question moves.

\smallskip
\begin{itemize}[itemsep=5pt]
  \item A rug's length is always \textbf{twice} its width. Then
        $A(w) = w \cdot 2w = $ \blank{1.6cm}, which is \blank{2.2cm},
        because \blank{1.2cm} lengths grow with the question.
  \item A rug's length is \textbf{fixed at 9 feet}. Then $A(w) = $
        \blank{1.6cm}, which is \blank{2.2cm}, because only
        \blank{1.2cm} length grows with the question.
\end{itemize}

\medskip
\textbf{Inverse proportion} means \TallMath{y = \frac{k}{x^n}} --- the question
sits \textit{underneath}, so as the question grows the answer
\blank{2.2cm}. These are \textbf{rational} functions; Unit 3 is about them, and
today you only have to recognize the words.
\end{notesbox}

\vspace{0.1in}

\boxguard[20]
\begin{notesbox}{4. Build It: the Constraint Answers for the Rest}
You picked a question. Every other quantity now has to become an
\textbf{answer} to it --- and the constraint is what forces them to.

\begin{enumerate}[itemsep=3pt]
  \item Name the \blank{2.2cm}, and name what you want out.
  \item Write the \blank{2.2cm}.
  \item Solve the constraint for the \textit{other} quantity.
  \item \blank{2.4cm} into the quantity you actually want.
\end{enumerate}

\medskip
\textbf{Ana's model.} Her question is the width $w$.

\smallskip
Step 2 --- the constraint: \quad \blank{1cm} $w \; + \;$ \blank{1cm} $= 80$

\smallskip
Step 3 --- solve it for the length: \quad $\ell = $ \blank{3.4cm}

\smallskip
Step 4 --- substitute into $A = w \cdot \ell$:

\begin{work}
  A(w) &= w \cdot \ell \\
       &= w(80 - 2w) \\
       &= 80w - 2w^2
\end{work}

\begin{center}
$A(w) = $ \blank{4.4cm}
\end{center}

Check it against Ana's table, where $A(20) = 800$:

\begin{work}
  A(20) &= 80(20) - 2(20)^2 \\
        &= 1600 - 800 \\
        &= 800
\end{work}

\smallskip
Part 3 predicted quadratic, and here is why it was right: the area multiplies
two lengths, and \textit{both} of them move with $w$ --- the length
$80 - 2w$ shrinks as the width grows.

\medskip
\textbf{Ben's model.} Same pen, same constraint, \textbf{different question}.
His question is the length $\ell$, so step 3 solves the same constraint for the
other letter.

\smallskip
Step 3: \quad $w = $ \blank{3.4cm}

\begin{work}
  A(\ell) &= w \cdot \ell \\
          &= \frac{80 - \ell}{2} \cdot \ell \\
          &= 40\ell - \frac{\ell^2}{2}
\end{work}

\begin{center}
$A(\ell) = $ \blank{4.4cm}
\end{center}

Check it against Ben's table, where $A(40) = 800$:

\begin{work}
  A(40) &= 40(40) - \frac{40^2}{2} \\
        &= 1600 - 800 \\
        &= 800
\end{work}

\medskip
\textbf{Two formulas. One pen.} Neither notebook is wrong. They differ because
the two members picked different \blank{3cm}.
\end{notesbox}

\vspace{0.1in}

\boxguard[20]
\begin{notesbox}{5. Use the Model}
A model that is never used is just algebra. Ana's is $A(w) = 80w - 2w^2$.

\medskip
\textbf{(a) Evaluate.}

\begin{work}
  A(10) &= 80(10) - 2(10)^2 \\
        &= 800 - 200 \\
        &= 600 \\
  A(30) &= 80(30) - 2(30)^2 \\
        &= 2400 - 1800 \\
        &= 600
\end{work}

$A(10) = $ \blank{1.8cm} ft$^2$ \qquad $A(30) = $ \blank{1.8cm} ft$^2$

\medskip
\textbf{(b) Average rate of change --- and the units name the question.}

\begin{work}
  \frac{A(20) - A(10)}{20 - 10} &= \frac{800 - 600}{10} \\
                                &= 20 \\
  \frac{A(30) - A(20)}{30 - 20} &= \frac{600 - 800}{10} \\
                                &= -20
\end{work}

On $[10, 20]$: \blank{1.6cm} \qquad On $[20, 30]$: \blank{1.6cm}

\smallskip
Units: \blank{6.4cm}

\smallskip
Straddle the best pen on \textbf{Ben's} model instead --- on $[30,40]$ and
$[40,50]$ his rates come out $+5$ and $-5$, in square feet per foot of
\blank{2cm}. Same pen, different numbers, because the units carry the question
you chose.

\medskip
\textbf{(c) The sign change locates the best pen.} The rate went from
\blank{1cm} to \blank{1cm}, so the largest area sits \textbf{between}. For
$A(w) = aw^2 + bw$ the vertex is at $w = -\dfrac{b}{2a}$:

\begin{work}
  w &= -\frac{b}{2a} \\
    &= -\frac{80}{2(-2)} \\
    &= 20
\end{work}

Best width $w = $ \blank{1.6cm} ft, \quad length $\ell = 80 - 2(20) = $
\blank{1.6cm} ft, \quad area $= $ \blank{2cm} ft$^2$.

\smallskip
\textbf{Back to the vote.} The best pen is \blank{1.4cm} ft by \blank{1.4cm} ft
--- \textbf{not} a square, because one side is free. Ben reaches the same pen
from the other end: his best is $\ell = $ \blank{1.6cm}.
\end{notesbox}

\vspace{0.1in}

\boxguard[20]
\begin{notesbox}{6. The Fine Print: Which Questions Are Allowed}
The algebra will happily answer questions the world cannot ask. The
\textbf{domain} is the set of questions the \textit{situation} allows.

\medskip
\begin{itemize}[itemsep=5pt]
  \item Ana's model: \blank{1.2cm} $< w <$ \blank{1.2cm}. \quad Why that upper
        end? At that width there is \blank{4.6cm}.
  \item Ben's model: \blank{1.2cm} $< \ell <$ \blank{1.2cm}.
  \item Same pen, \textbf{different domains} --- because the two members are
        asking different \blank{2.4cm}.
  \item Range: the area satisfies \blank{1.2cm} $< A \le$ \blank{1.6cm}
        square feet.
\end{itemize}

\smallskip
$A(-3) = -258$ is not bad arithmetic. It is a perfectly good answer to a
question the situation \blank{3.4cm}.

\medskip
\textbf{The assumption sentence.} Fill in the frame:

\smallskip
\begin{center}
\fbox{\begin{minipage}{0.92\linewidth}
\medskip
``This model assumes that \blank{4.4cm} , and it is only trustworthy for
inputs \blank{4cm} .''
\medskip
\end{minipage}}
\end{center}

\medskip
Name two things the pen model assumes about the real world:

\writeline

\medskip
\textbf{When one rule is not enough.} Suppose the barn wall is only 30 feet
long. Then a pen with $\ell > 30$ needs fencing on part of the fourth side too,
and the area rule \textbf{changes} at that threshold. A function with different
rules on different intervals is \blank{3cm} --- the last row of the table in
Part 3.
\end{notesbox}

\vspace{0.1in}

\boxguard
\begin{practicebox}
\begin{enumerate}[itemsep=8pt]
  \item The question is the \blank{1.2cm} row. Take differences until they are
        constant, then name the model type.

        \smallskip
        \renewcommand{\arraystretch}{1.4}
        \begin{center}
        \begin{tabular}{|c|c|c|c|c|c|}
        \hline
        $x$ & 0 & 1 & 2 & 3 & 4 \\
        \hline
        $y$ & 0 & 1 & 8 & 27 & 64 \\
        \hline
        1st & --- & \blank{1cm} & \blank{1cm} & \blank{1cm} & \blank{1cm} \\
        \hline
        2nd & --- & --- & \blank{1cm} & \blank{1cm} & \blank{1cm} \\
        \hline
        3rd & --- & --- & --- & \blank{1cm} & \blank{1cm} \\
        \hline
        \end{tabular}
        \end{center}

        Model type: \blank{3cm}

  \item A garden bed is built against a wall using 30 feet of edging on the
        other three sides. Take the \textbf{width} as your question, then build
        the area model and state its domain.

        \begin{work}
           2w + \ell &= 30 \\
                \ell &= 30 - 2w \\
                A(w) &= w(30 - 2w) \\
                     &= 30w - 2w^2
        \end{work}

        $A(w) = $ \blank{4cm} \qquad Domain: \blank{1cm} $< w <$ \blank{1cm}
\end{enumerate}
\end{practicebox}

\end{document}
